\bookchapter{Capítulo 5: Haz que el teléfono te entienda mejor — Ajustes básicos}

\chapterintro{Dedicar unos minutos a la configuración mejora enormemente la experiencia. Estos ajustes son como “afinar” el teléfono para que se ajuste a tus hábitos de uso.}

\booksection{I. Sonido y vibración}

Configura un tono de llamada y de notificación que te guste. Usa con flexibilidad el modo silencio (en reuniones) y el modo no molestar (al dormir); el mundo quedará en silencio al instante.

\begin{itemize}
  \item \textbf{Tono de llamada}: elige uno que puedas oír y que no dé vergüenza.
  \item \textbf{Modo vibración}: en silencio, avisa con la vibración.
  \item \textbf{Modo no molestar}: puedes definir excepciones, como permitir llamadas de contactos.
  \item \textbf{Volúmenes separados}: el volumen de multimedia, el de tono y el de alarma se pueden ajustar por separado.
\end{itemize}

\booksection{II. Pantalla y brillo}

\begin{itemize}
  \item \textbf{Brillo automático}: deja que el teléfono se ajuste según la luz ambiental, sin preocupaciones.
  \item \textbf{Modo protección ocular / Night Shift}: filtra la luz azul y hace que mirar el teléfono por la noche sea más cómodo.
  \item \textbf{Modo oscuro}: ahorra batería (especialmente eficaz en pantallas OLED) y por la noche no deslumbra. Muchas App también lo admiten.
  \item \textbf{Tasa de refresco de la pantalla}: algúnos Android permiten ajustarla; cuanto mayor (como 120 Hz), más fluido el desplazamiento.
  \item \textbf{Bloqueo automático}: ajústalo a 15 o 30 segundos; seguro y ahorra batería.
  \item \textbf{Tamaño de fuente}: ¿te parece pequeña? Auméntala.
\end{itemize}

\booksection{III. Fondo de pantalla, tema y escritorío}

Cambiar a un fondo bonito (estático o dinámico) cambia el ánimo. Los teléfonos Android de distintas marcas suelen tener una tienda de temas que permite cambiar todo el conjunto de iconos y estilo.

\textbf{Widgets de escritorío}: coloca el tiempo, la agenda o el control de música en el escritorío para acceder con un solo paso. iOS y Android lo admiten; mantén pulsado un espacio vacío del escritorío para añadirlos.

\booksection{IV. Ajustes de eficiencia: convierte el teléfono en tu “cerebro auxiliar”}

Esta parte reúne funciones “que una vez configuradas no hay vuelta atrás”; cada una puede mejorar notablemente tu eficiencia diaria.

\booksubsection{Modo concentración / No molestar: recupera la atención secuestrada por las notificaciones}

\begin{itemize}
  \item \textbf{Idea central}: no se trata de "mirar menos el teléfono", sino de "concentrarse cuando toca y comúnicarse cuando toque".
  \item \textbf{Recomendaciones de uso}:\newline · Modo trabajo: permite solo contactos "destacados" y App necesarias;\newline · Modo sueño: activación automática, permite solo la alarma;\newline · Ejercicio / lectura / reuniones: personalízalo según la actividad.
  \item \textbf{Usos avanzados}: con automatización + activación por ubicación/hora, los modos pueden conmutarse solos. Por ejemplo, en iOS Atajos puede configurar "al llegar a la oficina, entrar automáticamente en modo trabajo".
\end{itemize}

\booksubsection{Gestión de notificaciones: que el punto rojo deje de “acosarte”}

\begin{itemize}
  \item \textbf{Desactiva notificaciones de marketing}: el 90 \% de las notificaciones en realidad no las necesitas.
  \item \textbf{Ordena por importancia}: coloca en primer lugar las comunicaciones personales (teléfono, SMS, mensajería instantánea); las de noticias, compras y redes sociales pueden configurarse como "notificaciones silenciosas".
  \item \textbf{Agrupa notificaciones}: repliega las del mismo tipo para no quedar inundado por decenas de mensajes al desbloquear.
\end{itemize}

\booksubsection{Gestos rápidos: un paso por delante en la operación}

\begin{itemize}
  \item \textbf{Captura de pantalla}: botón de encendido + volumen abajo (iOS en pantalla completa) / botón de encendido + volumen abajo (general en Android).
  \item \textbf{Captura larga / captura con desplazamiento}: tras una captura normal, pulsa la opción "captura larga".
  \item \textbf{Escáner de código QR en un toque}: añade el icono "Escanear" al centro de control / ajustes rápidos.
  \item \textbf{Tocar atrás (iPhone)}: en [Configuración] > [Accesibilidad] > [Toque] > [Tocar atrás], asigna acciones como "captura de pantalla" o "bloquear".
  \item \textbf{Activación por voz}: Siri / Xiao Ai / Xiao Yi / Xiao Bu; con solo decirlo, puedes consultar el tiempo o poner una alarma.
\end{itemize}

\booksubsection{Teclado: la herramienta de eficiencia más infravalorada}

\begin{itemize}
  \item \textbf{Elige el teclado adecuado}:\newline · Sogou / Baidu: gran diccionarío, rápida actualización de neologismos;\newline · iFlytek: la mejor tasa de reconocimiento de voz, ideal para tomar notas en reuniones;\newline · Gboard (Google): mínimalista, fuerte sincronización multiplataforma, admite escritura deslizando;\newline · Teclado nativo de iOS: máxima privacidad, integración profunda con el sistema.
  \item \textbf{Interruptores imprescindibles}: corrección inteligente, predicción, sincronización de diccionarío en la nube, deslizar hacía arriba para escribir números y símbolos.
  \item \textbf{Advertencia de privacidad}: usa con precaución la "predicción de frases" o la "recomendación personalizada", podrían subir lo que escribes.
\end{itemize}

\booksubsection{Escritorío y multitarea: que el flujo de trabajo sea fluido}

\begin{itemize}
  \item \textbf{App en pantalla dividida / imagen en imagen}: responde mensajes mientras ves un video o toma notas mientras navegas.
  \item \textbf{Ventana pequeña / flotante}: WeChat / Douyin admiten "modo ventana pequeña"; al terminar, se recoge.
  \item \textbf{App frecuentes arriba}: en el Dock (parte inferior de la pantalla) coloca solo las 4 más usadas para reducir el tiempo de búsqueda.
\end{itemize}

\begin{tcbinfo}{Lista de comprobación del "primer día de trabajo"}
Al estrenar un teléfono o tras una gran actualización, dedica 10 minutos a repasar estos ajustes:\newline (1) Contraseña de bloqueo + biometría (2) Agrupar notificaciones + desactivar marketing (3) Tamaño de fuente / tasa de refresco / modo oscuro\newline (4) Teclado / navegador / correo predeterminados (5) Contacto de emergencia / tarjeta médica (6) Ajustes de copia de seguridad

\end{tcbinfo}

\booksection{V. Accesibilidad y adaptabilidad: el poder de la tecnología para el bien}

Que no te asuste el nombre “accesibilidad”; estas funciones son para todo el mundo. Se diseñaron inicialmente para personas con discapacidad, pero también son útiles para cualquier persona: con una mano lesionada o cargando al bebé, en el metro con los auriculares estropeados, con gafas de cerca leyendo textos pequeños... En definitiva, el diseño accesible encarna la idea de\textbf{la tecnología para el bien}: que todos puedan disfrutar por igual de la vida digital.

\booksubsection{Los “ojos” de los usuarios con discapacidad visual — guía completa del lector de pantalla}

El lector de pantalla es la herramienta central para que las personas con discapacidad visual usen el teléfono. Lee en voz alta\textbf{todo el contenido visible}de la pantalla; el usuario opera el teléfono con gestos y voz. A continuación se presentan los métodos de uso de VoiceOver en iOS y TalkBack en Android.

\begin{tcbinfo}{iOS VoiceOver}
\textbf{Ruta de activación}: [Configuración] > [Accesibilidad] > [VoiceOver], activa el interruptor. También puedes pulsar tres veces el botón lateral (o el botón de inicio) para activarlo/desactivarlo rápidamente.

\textbf{Gestos básicos}:

\begin{enumerate}
  \item Deslizar con un dedo: explorar el contenido de la pantalla elemento por elemento;
  \item Deslizar con dos dedos: desplazarse hacía arriba/abajo o izquierda/derecha por la página;
  \item Doble toque con un dedo: activar el elemento con el foco actual (equivalente a "clic");
  \item Deslizar con tres dedos: cambiar de página;
  \item Triple toque con un dedo: ampliar o seleccionar (distinto del triple toque normal).
\end{enumerate}

\textbf{Ajustes personalizados}: [VoiceOver] > [Voz], permite ajustar la\textbf{velocidad de lectura}(por defecto alrededor del 50 \%; se puede subir a 70-80 \% para una escucha más fluida), el\textbf{tono}y cambiar de paquete de voz. iOS incorpora “Voz mejorada” y “Controlador de voz”; el primero reduce respiraciones y pausas para una escucha más natural.

\end{tcbinfo}

\begin{tcbinfo}{Android TalkBack}
\textbf{Ruta de activación}: [Configuración] > [Accesibilidad] > [TalkBack], activa el interruptor. O mantén pulsadas las teclas de volumen arriba y abajo durante 3 segundos para activarlo rápidamente (algunas marcas lo admiten).

\textbf{Gestos básicos}:

\begin{enumerate}
  \item Deslizar a izquierda/derecha con un dedo: cambiar el elemento con el foco;
  \item Doble toque con un dedo: activar el elemento;
  \item Deslizar arriba/abajo o izquierda/derecha con dos dedos: desplazarse;
  \item Un toque con tres dedos: entrar en modo lectura (lectura continua);
  \item Un toque con cuatro dedos: pausar/reanudar la lectura.
\end{enumerate}

\textbf{Ajustes personalizados}: [TalkBack] > [Ajustes de texto a voz], permite ajustar la\textbf{velocidad},\textbf{tono}y cambiar de\textbf{motor de voz}(se recomienda instalar el motor de texto a voz de Google, con una lectura más natural).

\end{tcbinfo}

\booksubsection{De letras pequeñas a grandes, de borroso a nítido — ajustes visuales}

\booksubsection{Gesto de ampliación}

\begin{itemize}
  \item \textbf{iOS}: [Configuración] > [Accesibilidad] > [Zoom], al activarlo,\textbf{tres toques en la pantalla}amplían. Por defecto amplía toda la pantalla; puede elegirse ampliación en ventana (lupa parcial).\textbf{Se ajusta el nivel (1,5x a 15x) pellizcando con dos dedos}; un doble toque con tres dedos bloquea la ampliación actual.
  \item \textbf{Android}: [Configuración] > [Accesibilidad] > [Ampliación], tres toques para ampliar y otros tres para salir; o directamente\textbf{tres toques seguidos con tres dedos}para una ampliación temporal (hasta 8x).
\end{itemize}

\booksubsection{Texto de alto contraste / texto en negrita}

\begin{itemize}
  \item \textbf{iOS}: [Configuración] > [Accesibilidad] > [Visualización y tamaño del texto], activa "Texto en negrita", "Aumentar contraste", "Reducir transparencia", etc. "Inversión inteligente" mantiene el texto en negro sobre fondos claros y evita texto blanco sobre fondo blanco.
  \item \textbf{Android}: [Configuración] > [Pantalla] > [Tamaño de fuente], con cuatro niveles: "Pequeño, Normal, Grande, Extra grande". Huawei/Xiaomi añaden interruptores de "Negrita" y "Texto de alto contraste".
\end{itemize}

\booksubsection{Filtros de color (aptos para daltonismo)}

\begin{itemize}
  \item \textbf{iOS}: [Configuración] > [Accesibilidad] > [Visualización y tamaño del texto] > [Filtros de color], ofrece\textbf{cuatro modos para daltonismo rojo/verde/azul/amarillo}, además de "Escala de grises" e "Invertir colores". Se recomienda hacer antes un test de visión de colores (integrado en el sistema) y elegir el filtro correspondiente.
  \item \textbf{Android}: [Configuración] > [Accesibilidad] > [Corrección de color], con\textbf{tres modos para daltonismo}+ un interruptor de "Invertir colores".
\end{itemize}

\booksubsection{Descripción de imágenes con IA: que la IA sea tus “ojos”}

Cuando el lector de pantalla se topa con una imagen, lo habitual era saltarla o leer solo la palabra “imagen”. Pero hoy la IA puede\textbf{describir el contenido de la imagen en una frase}, muy útil para reconocer captchas, leer díagramas o identificar menús.

\begin{tcbinfo}{Soporte de descripción de imágenes con IA por sistema}
\begin{itemize}
  \item \textbf{iOS 15+}: en [Configuración] > [Accesibilidad] > [VoiceOver] > [Descripciones de voz], activa "Descripciones de imágenes". El sistema invoca un\textbf{modelo de IA local}para generar automáticamente descripciiones textuales (requiere iPhone 12 o superior).
  \item \textbf{TalkBack en versiones recientes}: Google TalkBack integra en sus versiones más recientes la función de\textbf{preguntas sobre imágenes}, que permite preguntar "¿qué hay en esta imagen?" y responder en tiempo real (no depende estrictamente de la versión de Android; basta con actualizar la App TalkBack).
  \item \textbf{App de terceros}: "Seeing AI" (iOS), "Lookout" (Android) y "Xiao Yi Bang Bang" (Huawei) son excelentes asistentes visuales con IA; pueden reconocer objetos, leer documentos e identificar expresiiones faciales.
\end{itemize}

\end{tcbinfo}

\begin{tcbpractice}{Escenarios de uso}
Prueba a preguntar sobre una foto de calle “¿cuántas personas hay en esta imagen? ¿Qué están haciendo?”, o abre una App de comida a domicilio y deja que la IA lea los nombres de los platos del menú en imagen —\textbf{esto es verdaderamente “la tecnología hace la vida más igualitaria”}.

\end{tcbpractice}

\booksubsection{Los “oídos” de los usuarios con discapacidad auditiva — guía completa de subtítulos en tiempo real}

Los subtítulos en tiempo real (Live Captíons) son una función\textbf{a nivel de sistema que convierte voz a texto}; son especialmente útiles para personas con discapacidad auditiva en entornos ruidosos, bibliotecas o situaciones en otros idiomas.

\begin{tcbinfo}{Activación y optimización de subtítulos en tiempo real}
\textbf{iOS}: [Configuración] > [Accesibilidad] > [Subtítulos en tiempo real], activa el interruptor. El sistema detecta todo el audio del dispositivo (llamadas, videos, podcasts, FaceTime, notas de voz) y muestra los subtítulos automáticamente en la parte superior de la pantalla.

\textbf{Android 12+}: pulsa la\textbf{tecla de volumen}toca el icono “Subtítulos en tiempo real” en la ventana emergente; o actívalo en [Configuración] > [Accesibilidad] > [Subtítulos en tiempo real].

\textbf{Consejos para mejorar la precisión}:

\begin{enumerate}
  \item Úsalo preferentemente en un\textbf{entorno silencioso}, lejos de aires acondiciónados o ruidos de máquinas;
  \item La precisión es mayor cuando el hablante\textbf{articula con claridad y a velocidad moderada};
  \item Las conversaciones entre varias personas reducen la precisión; puede usarse el modo "mostrar solo a una persona";
  \item Admite mandarín, cantonés, inglés y otros idiomas principales;\textbf{la precisión en idiomas minoritaríos es menor};
  \item Si hay música de fondo o voces que interfieren durante la conversación, el reconocimiento se interrumpirá.
\end{enumerate}

\end{tcbinfo}

\booksubsection{Escenarío de llamada bidireccional: voz a texto + texto a voz}

Cuando el usuario\textbf{ni puede oír ni puede hablar}(afonía súbita, recuperación de cirugía de laringe, tartamudez severa), el teléfono puede hacer de “intérprete”.

\begin{tcbinfo}{Transcripción en tiempo real de llamadas (TTY / RTT)}
\textbf{iOS}: [Configuración] > [Accesibilidad] > [RTT/TTY], al activarlo, durante la llamada aparece un teclado en pantalla; la voz del interlocutor se transcribe a texto en tiempo real y\textbf{lo que tú escribas se lee en voz alta}al interlocutor.

\textbf{Android}: [Configuración] > [Accesibilidad] > [Transcripción en tiempo real] + [Avisos de sonido]. “Transcripción en tiempo real” muestra lo que dice el interlocutor como texto; “Avisos de sonido” convierte lo que tú escribas en voz cuando no es posible hablar.

\textbf{Compatibilidad}:

\begin{itemize}
  \item \textbf{RTT (Real-Time Text)}: ambos interlocutores deben ser compatibles; el texto se muestra a medida que se escribe, sin latencia; los principales operadores ya lo admiten;
  \item \textbf{TTY (teletipo)}: requiere un dispositivo TTY específico; está siendo sustituido por RTT;
  \item La\textbf{función de voz a texto}de WeChat/QQ es una alternativa: mantén pulsado el mensaje de voz "Convertir a texto".
\end{itemize}

\end{tcbinfo}

\booksubsection{Funciones “pequeñas pero útiles” para todos}

\begin{itemize}
  \item \textbf{Tocar atrás}(iOS 14+): [Configuración] > [Accesibilidad] > [Toque] > [Tocar atrás], permite asignar a\textbf{doble/triple toque}acciones como captura de pantalla, bloqueo, centro de control, etc. Muy útil cuando llevas paquetes con ambas manos.
  \item \textbf{AssistiveTouch / botón flotante}: en iOS [Accesibilidad] > [Toque] > [AssistiveTouch]; en Android [Accesibilidad] > [Navegación flotante]. Un botón virtual que sustituye gestos o botónes físicos; muy útil cuando un botón falla o para usar con una mano.
  \item \textbf{Control por interruptores}: controla el teléfono mediante interruptores externos, toques únicos o movimientos de la cabeza,\textbf{diseñado para personas con discapacidad motora}. En iOS [Accesibilidad] > [Control por interruptores]; en Android [Accesibilidad] > [Control por interruptores].
  \item \textbf{Reconocimiento de sonidos}(iOS): [Accesibilidad] > [Reconocimiento de sonidos], al activarlo, el teléfono "escucha" timbres, llantos de bebé, alarmás de incendio, sirenas, etc., y\textbf{al detectarlos avisa con vibración + notificación en pantalla}; muy valioso para usuarios con discapacidad auditiva y durante el sueño.
  \item \textbf{Aviso parpadeante LED}: en iOS [Accesibilidad] > [Audio/Visual] > [Aviso LED parpadeante]; en Android [Configuración] > [Notificaciones] > [Aviso por flash]. El flash parpadea al recibir llamadas y notificaciones;\textbf{incluso en entornos ruidosos se puede percibir sin sonido}.
  \item \textbf{Audio mono}: [Configuración] > [Accesibilidad] > [Audio] > [Audio mono], combina los canales izquierdo y derecho. Los usuarios con pérdida auditiva de un oído no pierden la mitad de la información sonora.
\end{itemize}

\begin{tcbinfo}{Filosofía de accesibilidad de Apple}
Apple invierte mucho en accesibilidad y cada año actualiza y añade funciones. La web oficial de Apple tiene una página específica de accesibilidad que detalla el uso de cada función. Las funciones de accesibilidad de iOS suelen ser más completas y fluidas que las de Android, una razón importante por la que muchos usuarios con discapacidad eligen Apple.

\end{tcbinfo}

\begin{tcbpractice}{Dedica 30 segundos a probarlo}
Busca un rato libre y haz estas pruebas:

\begin{enumerate}
  \item Activa\textbf{VoiceOver / TalkBack}e intenta navegar por WeChat deslizando con tres dedos sin usar los ojos;
  \item Activa los\textbf{subtítulos en tiempo real}y habla al teléfono para comprobar la precisión;
  \item Usa\textbf{tres toques}para ampliar un texto al 200 \% y experimentar el día a día de una persona con discapacidad visual;
  \item Desactiva todas las funciones anteriores y reflexiona: "Si siempre necesitara estas funciones, ¿cómo querría que estuvieran diseñadas?".
\end{enumerate}

\end{tcbpractice}
